\documentclass[11pt,a4paper]{article}
\usepackage{shared/parscover}

% ============================================================================
% Coercion-Resistant Protocols for Governance Under Authoritarian Surveillance
% Pars Network Technical Whitepaper PNP-005
% ============================================================================

\usepackage[utf8]{inputenc}
\usepackage[T1]{fontenc}
\usepackage{amsmath,amssymb,amsthm}
\usepackage{graphicx}
\usepackage{hyperref}
\usepackage{algorithm}
\usepackage{algpseudocode}
\usepackage{booktabs}
\usepackage{listings}
\usepackage{xcolor}
\usepackage{tikz}
\usepackage{enumitem}
\usepackage{caption}
\usepackage{fancyhdr}
\usepackage{abstract}

\geometry{margin=1in}

\hypersetup{
    colorlinks=true,
    linkcolor=blue!70!black,
    citecolor=green!50!black,
    urlcolor=blue!60!black
}

\lstset{
    basicstyle=\ttfamily\small,
    breaklines=true,
    frame=single,
    backgroundcolor=\color{gray!10},
    keywordstyle=\color{blue!70!black}\bfseries,
    commentstyle=\color{green!50!black}\itshape,
    stringstyle=\color{red!60!black},
    numbers=left,
    numberstyle=\tiny\color{gray},
    numbersep=5pt
}

\usetikzlibrary{arrows.meta,positioning,shapes.geometric,fit,calc}

\newtheorem{definition}{Definition}[section]
\newtheorem{theorem}{Theorem}[section]
\newtheorem{lemma}{Lemma}[section]
\newtheorem{proposition}{Proposition}[section]
\newtheorem{corollary}{Corollary}[section]

\pagestyle{fancy}
\fancyhf{}
\fancyhead[L]{\small Cyrus Pahlavi, Pars Team}
\fancyhead[R]{\small Coercion Resistance}
\fancyfoot[C]{\thepage}

% ============================================================================
\title{%
    \textbf{Coercion-Resistant Protocols for Governance\\
    Under Authoritarian Surveillance}\\[0.5em]
    \large Pars Network Technical Whitepaper PNP-005
}

\author{%
    Cyrus Pahlavi, Pars Team\\
    \texttt{research@pars.network}\\[0.5em]
    \small Version 1.0 --- February 2026
}

\date{}

% ============================================================================
\begin{document}
\parscoverpage

\thispagestyle{fancy}

% ============================================================================
\begin{abstract}
\noindent
We present a suite of coercion-resistant protocols for the Pars Network,
designed to protect diaspora community members who participate in
decentralised governance under threat of physical coercion by authoritarian
regimes.  Standard cryptographic protocols assume users freely control their
keys; this assumption fails under detention, threat, or violence.  We
introduce four mechanisms: (1)~deniable authentication producing transcripts
indistinguishable from innocent interactions; (2)~panic passwords triggering
credential rotation and evidence destruction while presenting a plausible
decoy environment; (3)~decoy wallets with credible on-chain histories; and
(4)~plausible deniability for governance votes through re-randomisable
ballot encryption.  We formalise coercion resistance in the Universal
Composability framework, prove receipt-freeness and coercion resistance
under standard assumptions, and present a threat-scenario analysis based
on documented practices of surveillance states.  The protocols integrate
with the Pars session layer (PNP-003) and governance system (PNP-001).
\end{abstract}

\vspace{1em}
\noindent\textbf{Keywords:} coercion resistance, deniable authentication,
panic passwords, decoy wallets, plausible deniability, governance privacy

\tableofcontents
\newpage

% ============================================================================
\section{Introduction}
\label{sec:introduction}

Participation in decentralised governance is a political act.  For the
Persian diaspora, especially those with family inside Iran, governance
participation carries material risk.  The Islamic Republic has a documented
history of coercing individuals to reveal digital credentials, confiscating
devices, and retaliating against diaspora
activists' families~\cite{freedomhouse2023iran}.

Standard cryptographic threat models assume private keys are under sole
owner control.  This collapses under physical coercion.

\subsection{Coercion Scenarios}

\begin{enumerate}[label=\textbf{S\arabic*:}]
    \item \textbf{Border Crossing:} Device inspected; user compelled to
    unlock and demonstrate contents.
    \item \textbf{Detention:} Interrogation demanding passwords, wallet
    keys, communication histories.
    \item \textbf{Family Pressure:} Threats against family members to
    compel vote record disclosure.
    \item \textbf{Warrant:} Cooperative jurisdiction compels provider to
    reveal participation records.
\end{enumerate}

\subsection{Design Goals}

\begin{enumerate}
    \item \textbf{Coercion Resistance:} Coerced user complies without
    revealing true governance participation or holdings.
    \item \textbf{Receipt-Freeness:} No user can prove how they voted.
    \item \textbf{Plausible Deniability:} Existence of true credentials
    indistinguishable from non-existence.
    \item \textbf{Minimal Overhead:} $<2\times$ overhead over normal
    operations.
\end{enumerate}

\subsection{Contributions}

\begin{itemize}
    \item Formalisation of coercion resistance for blockchain governance.
    \item Deniable authentication with offline deniability.
    \item Panic password system with cryptographic decoy generation.
    \item Decoy wallets with credible transaction histories.
    \item Re-randomisable ballot encryption for receipt-free voting.
    \item Threat-scenario analysis from documented state practices.
\end{itemize}

% ============================================================================
\section{Threat Model}
\label{sec:threat}

\begin{definition}[Coercive Adversary]
    $\mathcal{A}_{\mathrm{coerce}}$ has network adversary capabilities
    and can additionally:
    \begin{enumerate}[label=(\alph*)]
        \item Compel secret disclosure
        \item Compel any protocol action
        \item Inspect device state
        \item Observe physical behaviour during execution
    \end{enumerate}
\end{definition}

\begin{definition}[Coercion Resistance]
    Protocol $\Pi$ is coercion-resistant if strategy
    $\sigma_{\mathrm{fake}}$ exists such that a coerced user following it
    produces interaction indistinguishable from any $\sigma'$, from
    $\mathcal{A}_{\mathrm{coerce}}$'s perspective.
\end{definition}

\begin{definition}[Receipt-Freeness]
    A voting protocol is receipt-free if no voter can construct a receipt
    proving their vote, even cooperating with the receipt requester.
\end{definition}

% ============================================================================
\section{Deniable Authentication}
\label{sec:deniable-auth}

\subsection{Problem}

Digital signatures produce non-repudiable transcripts.  If coerced to
reveal authentication history, signatures prove specific interactions.

\subsection{Ring-Signature-Based Authentication}

\begin{algorithm}[H]
\caption{Deniable Authentication}
\label{alg:deniable-auth}
\begin{algorithmic}[1]
\Require Prover $A$ with $(a, A)$, Verifier $B$ with $(b, B)$
\State $A$ generates ephemeral $(r, R = rG)$
\State $\sigma \gets \mathrm{RingSign}(\{A, B\}, a, m)$
\State $A \to B$: $(m, R, \sigma)$
\State $B$ verifies: $\mathrm{RingVerify}(\{A, B\}, m, \sigma)$
\State Third party cannot tell whether $A$ or $B$ signed
\end{algorithmic}
\end{algorithm}

\begin{theorem}[Deniability]
    Ring-signature authentication produces simulatable transcripts.
    No third party can determine which ring member signed.
\end{theorem}

\begin{proof}
    By ring signature anonymity~\cite{rivest2001leak}, $\sigma$ from $A$
    is computationally indistinguishable from $\sigma$ from $B$.  Verifier
    $B$ can produce valid ring signatures using their own key, simulating
    any transcript.
\end{proof}

\subsection{PSP Integration}

PSP (PNP-003) already provides pairwise deniability through X3DH's use
of DH.  Ring-signature authentication extends this to explicit
authentication contexts (governance proposals, identity assertions).

% ============================================================================
\section{Panic Passwords}
\label{sec:panic}

\subsection{Concept}

\begin{definition}[Panic Password System]
    \begin{itemize}
        \item Primary password $p_1$: unlocks true environment
        \item Panic password $p_2$: unlocks decoy, triggers protection
        \item $\mathrm{KDF}(p_1) \neq \mathrm{KDF}(p_2)$; both produce
        valid-looking keys
    \end{itemize}
\end{definition}

\subsection{Dual-Key Derivation}

\begin{algorithm}[H]
\caption{Dual-Key Derivation}
\label{alg:dual-key}
\begin{algorithmic}[1]
\Procedure{Setup}{$p_1$, $p_2$}
    \State $k_1 \gets \mathrm{Argon2id}(p_1, s_1)$
    \State $k_2 \gets \mathrm{Argon2id}(p_2, s_2)$
    \State $E_1 \gets \mathrm{AES}(k_1, \mathrm{state}_{\mathrm{true}})$
    \State $E_2 \gets \mathrm{AES}(k_2, \mathrm{state}_{\mathrm{decoy}})$
    \State Store $(E_1, E_2, s_1, s_2)$
\EndProcedure
\Procedure{Unlock}{$p$}
    \State Try $k \gets \mathrm{Argon2id}(p, s_1)$;
        if $\mathrm{AES.Dec}(k, E_1)$ succeeds: return true env
    \State Try $k \gets \mathrm{Argon2id}(p, s_2)$;
        if $\mathrm{AES.Dec}(k, E_2)$ succeeds:
        \Statex \quad TriggerPanic(); return decoy env
    \State Return failure
\EndProcedure
\end{algorithmic}
\end{algorithm}

\subsection{Panic Actions}

\begin{enumerate}
    \item \textbf{Key Rotation:} Rotate session keys and prekey bundles.
    \item \textbf{Evidence Destruction:} Secure-erase primary state $E_1$.
    \item \textbf{Silent Alert:} Dead-man notification via mesh to
    designated contacts.
    \item \textbf{Decoy Presentation:} Show plausible decoy environment.
\end{enumerate}

\begin{lstlisting}[language=Go,caption={Panic Action Handler}]
type PanicHandler struct {
    keys     *KeyManager
    storage  *SecureStorage
    mesh     *MeshClient
    contacts []PeerID
}

func (h *PanicHandler) Execute(
    ctx context.Context,
) error {
    // Rotate keys (best-effort)
    h.keys.RotateAll(ctx)

    // Destroy primary state
    if err := h.storage.SecureErase(
        "primary_state",
    ); err != nil {
        return fmt.Errorf("erase: %w", err)
    }

    // Silent alert via mesh
    alert := &CoercionAlert{
        Time: time.Now().Unix(),
        Type: AlertPanic,
    }
    for _, c := range h.contacts {
        h.mesh.SendAsync(ctx, c, alert)
    }
    return nil
}
\end{lstlisting}

\begin{proposition}[Post-Panic Indistinguishability]
    After panic activation, an inspector cannot determine whether a
    primary state ever existed.
\end{proposition}

\begin{proof}[Proof sketch]
    After secure erasure of $E_1$, only $(E_2, s_1, s_2)$ remains.
    $s_1$ and zeroed $E_1$ space are indistinguishable from unused
    storage.  The device appears to contain a single encrypted state
    ($E_2$), unlockable with the panic password presented as
    ``the only'' password.
\end{proof}

% ============================================================================
\section{Decoy Wallets}
\label{sec:decoy}

\subsection{Requirements}

\begin{enumerate}
    \item \textbf{Credibility:} Plausible balance and history.
    \item \textbf{Verifiability:} Real on-chain transactions.
    \item \textbf{Consistency:} Matches claimed usage pattern.
    \item \textbf{Independence:} No on-chain link to true wallet.
\end{enumerate}

\subsection{Construction}

\begin{algorithm}[H]
\caption{Decoy Wallet Construction}
\label{alg:decoy-wallet}
\begin{algorithmic}[1]
\Procedure{CreateDecoy}{profile $P$}
    \State $(d, D) \gets$ fresh key pair
    \State Fund via stealth address from decoy fund
    \For{$i = 1$ to $P.\mathrm{txCount}$}
        \State Schedule tx at $t_{\mathrm{now}} - P.\mathrm{spacing} \cdot i$
        \State Amount $\sim P.\mathrm{amountDist}$
        \State Recipient $\in P.\mathrm{counterparties}$
    \EndFor
    \State Maintain $P.\mathrm{targetBalance}$
    \State Store under panic key $k_2$
\EndProcedure
\end{algorithmic}
\end{algorithm}

\begin{lstlisting}[language=Go,caption={Decoy Profile}]
type DecoyProfile struct {
    MinBalance   *big.Int
    MaxBalance   *big.Int
    TxPerWeek    int
    AmountMean   *big.Int
    AmountStdDev *big.Int
    Counterparties []common.Address
    Tokens       []TokenHolding
    VoteFreq     float64
}

var Standard = DecoyProfile{
    MinBalance: eth(0.1),
    MaxBalance: eth(2.0),
    TxPerWeek:  3,
    AmountMean: eth(0.05),
    AmountStdDev: eth(0.03),
    VoteFreq:   0.3,
}
\end{lstlisting}

\subsection{Funding Unlinkability}

\begin{enumerate}
    \item Decoy funds from community-managed pool with ring-signature
    deposits~\cite{noether2015ring}.
    \item Transfers via stealth addresses (PNP-007).
    \item Pool replenished via constant-rate treasury drip.
\end{enumerate}

% ============================================================================
\section{Receipt-Free Governance Voting}
\label{sec:voting}

\subsection{Re-Randomisable Ballots}

\begin{definition}[Re-Randomisable Ballot]
    Ballot $B = (g^r,\; v \cdot h^r)$ under ElGamal with PK $h = g^x$
    is re-randomised:
    \[
        B' = (g^{r+r'},\; v \cdot h^{r+r'})
    \]
    by any party knowing $h$, using fresh $r'$.
\end{definition}

\subsection{Voting Protocol}

\begin{algorithm}[H]
\caption{Coercion-Resistant Vote Casting}
\label{alg:cr-vote}
\begin{algorithmic}[1]
\Require Voter $V$, vote $v$, election PK $h$
\State \textbf{Registration (anonymous):}
\State $V$ obtains credential $\mathrm{cred}$ via ZK eligibility proof
\State $V$ also registers fake credential $\mathrm{cred}_{\mathrm{fake}}$
\State
\State \textbf{Casting:}
\State $B \gets \mathrm{ElGamal.Enc}(h, v, r)$
\State $\pi \gets \mathrm{Prove}(B \in \{\mathrm{Enc}(h,0), \mathrm{Enc}(h,1)\})$
\State Submit $(B, \pi, \mathrm{cred})$ via anonymous channel
\State
\State \textbf{Re-randomisation (trustees):}
\ForAll{trustee $T_j$}
    \State $B \gets \mathrm{ReRand}(B, r'_j)$
\EndFor
\State After this, $V$ cannot link to final ballot
\State
\State \textbf{Under coercion:} cast with $\mathrm{cred}_{\mathrm{fake}}$
\State Tally system silently ignores fake-credential votes
\end{algorithmic}
\end{algorithm}

\begin{theorem}[Receipt-Freeness]
    After trustee re-randomisation, no voter can construct a receipt
    linking themselves to a re-randomised ballot.
\end{theorem}

\begin{proof}
    Original ballot $B = (g^r, v \cdot h^r)$ becomes
    $B' = (g^{r+r'}, v \cdot h^{r+r'})$ with $r' = \sum_j r'_j$ unknown
    to the voter.  Even revealing original $r$, the voter cannot link $B$
    to $B'$.  By ElGamal IND-CPA security, $B'$ is indistinguishable from
    any ballot for the same vote value.
\end{proof}

\begin{theorem}[Coercion Resistance]
    A coerced voter using the fake-credential strategy is
    indistinguishable from honest voting.
\end{theorem}

\begin{proof}
    $\mathrm{cred}_{\mathrm{fake}}$ is syntactically identical to
    $\mathrm{cred}$ (both are ZK proofs).  The coercer cannot distinguish
    them without credential issuer access.  The vote with
    $\mathrm{cred}_{\mathrm{fake}}$ is encrypted identically.  The tally
    system silently rejects it; the coercer cannot detect rejection.
\end{proof}

\subsection{Trustee Threshold}

$t$-of-$n$ threshold with recommended $t = 5$, $n = 9$, trustees from
distinct jurisdictions.

\begin{definition}[Threshold Re-Randomisation]
    Election PK $h = g^x$ with $x = \sum_{j=1}^n x_j \bmod q$ shared
    among $n$ trustees.  Re-randomisation requires $\geq t$ contributors.
\end{definition}

% ============================================================================
\section{Plausible Deniability for Governance}
\label{sec:deniability}

\subsection{Vote Deniability}

\begin{enumerate}
    \item \textbf{Credential Hiding:} Encrypted under primary key;
    destroyed on panic.
    \item \textbf{No On-Chain Identity:} Votes via session layer,
    aggregated by trustees before L2 commitment.
    \item \textbf{Time-Bounded Records:} Encryption keys rotated per
    epoch; individual data deleted after tally.
\end{enumerate}

\begin{proposition}[Participation Deniability]
    No adversary can determine with non-negligible advantage whether a
    specific user participated.
\end{proposition}

\begin{proof}[Proof sketch]
    Ballots submitted via anonymous routing, re-randomised by trustees.
    No on-chain voter-ballot mapping.  Device after panic contains no
    credential evidence.
\end{proof}

% ============================================================================
\section{Dead Man's Switch}
\label{sec:dead-man}

\begin{algorithm}[H]
\caption{Dead Man's Switch}
\label{alg:dead-man}
\begin{algorithmic}[1]
\Procedure{Setup}{user $U$, interval $T$, contacts $C$}
    \State Register heartbeat schedule with interval $T$
    \State Pre-encrypt alert messages for each $c \in C$
    \State Distribute time-locked capsules to relay network
\EndProcedure
\Procedure{Heartbeat}{}
    \State Authenticated heartbeat to relay; reset timer
\EndProcedure
\Procedure{Trigger}{}
    \State Timer expires without heartbeat
    \State Decrypt and deliver alerts; rotate $U$'s keys;
    publish DID revocation
\EndProcedure
\end{algorithmic}
\end{algorithm}

\begin{definition}[Time-Lock Capsule]
    $C = (N, a, t, Z)$ where $N = pq$, $Z = m \oplus H(a^{2^t} \bmod N)$.
    Openable in $O(t)$ sequential squarings; not
    parallelisable~\cite{rivest1996time}.
\end{definition}

% ============================================================================
\section{Threat Scenario Analysis}
\label{sec:analysis}

\begin{table}[h]
\centering
\caption{Mitigation Matrix}
\label{tab:mitigation}
\begin{tabular}{@{}lll@{}}
\toprule
\textbf{Scenario} & \textbf{Threat} & \textbf{Mitigation} \\
\midrule
S1: Border    & Device inspection    & Decoy wallet + panic pw \\
S1: Border    & Password demand      & Panic pw as primary \\
S2: Detention & Extended interrogation & Panic destroys evidence \\
S2: Detention & Key disclosure       & Decoy keys presented \\
S2: Detention & Missed check-in      & Dead man's switch \\
S3: Family    & Vote disclosure      & Receipt-free voting \\
S3: Family    & Coerced voting       & Fake credentials \\
S4: Warrant   & Server data request  & No plaintext server-side \\
S4: Warrant   & Blockchain analysis  & Stealth addresses \\
\bottomrule
\end{tabular}
\end{table}

% ============================================================================
\section{Implementation}
\label{sec:implementation}

\subsection{Secure Erasure}

\begin{lstlisting}[language=Go,caption={Secure Erasure}]
func (s *SecureStorage) Erase(
    key string,
) error {
    path := filepath.Join(s.path, key)
    info, err := os.Stat(path)
    if err != nil {
        return fmt.Errorf("stat: %w", err)
    }
    size := info.Size()
    f, err := os.OpenFile(
        path, os.O_WRONLY, 0,
    )
    if err != nil {
        return fmt.Errorf("open: %w", err)
    }
    defer f.Close()

    buf := make([]byte, 4096)
    for pass := 0; pass < 3; pass++ {
        f.Seek(0, 0)
        rem := size
        for rem > 0 {
            n := min(int64(len(buf)), rem)
            rand.Read(buf[:n])
            f.Write(buf[:n])
            rem -= n
        }
        f.Sync()
    }
    f.Close()
    return os.Remove(path)
}
\end{lstlisting}

\subsection{App Disguise}

\begin{enumerate}
    \item Icon and name changeable to common utility (calculator, weather).
    \item Disguise activated via specific gesture.
    \item Functional utility interface; Pars UI behind secondary gesture.
\end{enumerate}

% ============================================================================
\section{Formal Security Analysis}
\label{sec:formal}

\begin{definition}[$\mathcal{F}_{\mathrm{CR}}$]
    Ideal functionality accepting votes from registered voters.  On
    coercion: accepts ``coerced vote'' without counting it, reveals same
    view as honest vote to coercer, tallies only non-coerced votes.
\end{definition}

\begin{theorem}[UC Realisation]
    The Pars coercion-resistant voting protocol UC-realises
    $\mathcal{F}_{\mathrm{CR}}$ under DDH, assuming $< t$ corrupted
    trustees.
\end{theorem}

The proof follows Juels et al.~\cite{juels2005coercion}, extended with
anonymous credential and session-layer integration.

% ============================================================================
\section{Limitations}
\label{sec:limitations}

\begin{enumerate}
    \item No protocol defeats sustained physical torture; operational
    security (ephemeral devices) is necessary for high-risk situations.
    \item Decoy credibility requires regular profile updates.
    \item If adversary knows panic system exists, they may demand both
    passwords.
    \item $\geq t$ colluding trustees can de-anonymise votes.
\end{enumerate}

% ============================================================================
\section{Related Work}
\label{sec:related}

Juels et al.~\cite{juels2005coercion} introduced coercion-resistant
e-voting.  Civitas~\cite{clarkson2008civitas} implemented the scheme.
Chaum~\cite{chaum2004receipt} provides receipt-free verification.
TrueCrypt/VeraCrypt~\cite{truecrypt2004,veracrypt2015} pioneered
plausible-deniability disk encryption.  Briar~\cite{briar2017} has a
panic button.  Pars integrates coercion resistance into a blockchain
governance platform with diaspora-specific threat modelling.

% ============================================================================
\section{Conclusion}
\label{sec:conclusion}

The Pars coercion-resistant suite addresses physical security threats
facing diaspora members.  Deniable authentication, panic passwords, decoy
wallets, and receipt-free voting provide layered defence.  Formal analysis
in the UC framework proves coercion resistance and receipt-freeness under
standard assumptions.

Cryptographic security alone is insufficient when adversaries operate
outside standard threat model assumptions.  Explicitly modelling coercion
as a protocol-level threat enables meaningful protection even when physical
safety is compromised.

% ============================================================================
\bibliographystyle{plain}
\begin{thebibliography}{26}

\bibitem{freedomhouse2023iran}
Freedom House,
``Freedom on the Net: Iran,''
2023.

\bibitem{rivest2001leak}
R.~Rivest, A.~Shamir, and Y.~Tauman,
``How to Leak a Secret,''
\emph{ASIACRYPT 2001}, pp.~552--565, 2001.

\bibitem{unger2015dake}
N.~Unger and I.~Goldberg,
``Deniable Key Exchanges for Secure Messaging,''
\emph{CCS}, pp.~1211--1223, 2015.

\bibitem{juels2005coercion}
A.~Juels, D.~Catalano, and M.~Jakobsson,
``Coercion-Resistant Electronic Elections,''
\emph{WPES}, pp.~61--70, 2005.

\bibitem{clarkson2008civitas}
M.~Clarkson, S.~Chong, and A.~Myers,
``Civitas: Toward a Secure Voting System,''
\emph{IEEE S\&P}, pp.~354--368, 2008.

\bibitem{chaum2004receipt}
D.~Chaum,
``Secret-ballot receipts,''
\emph{IEEE Security \& Privacy}, 2(1), 2004.

\bibitem{truecrypt2004}
TrueCrypt Foundation,
``TrueCrypt: Hidden Volumes,''
2004.

\bibitem{veracrypt2015}
IDRIX,
``VeraCrypt: Plausible Deniability,''
2015.

\bibitem{briar2017}
Briar Project,
``Briar: Secure messaging, anywhere,''
2017.

\bibitem{noether2015ring}
S.~Noether,
``Ring Signature Confidential Transactions for Monero,''
\emph{MRL-0005}, 2015.

\bibitem{rivest1996time}
R.~Rivest, A.~Shamir, and D.~Wagner,
``Time-lock puzzles and timed-release crypto,''
\emph{MIT/LCS/TR-684}, 1996.

\bibitem{canetti2001universally}
R.~Canetti,
``Universally Composable Security,''
\emph{FOCS}, pp.~136--145, 2001.

\bibitem{goldwasser1989knowledge}
S.~Goldwasser, S.~Micali, and C.~Rackoff,
``The Knowledge Complexity of Interactive Proof Systems,''
\emph{SIAM J.\ Comput.}, 18(1), 1989.

\bibitem{chaum1981untraceable}
D.~Chaum,
``Untraceable electronic mail,''
\emph{Commun.\ ACM}, 24(2), 1981.

\bibitem{pedersen1991threshold}
T.~Pedersen,
``A Threshold Cryptosystem without a Trusted Party,''
\emph{EUROCRYPT '91}, pp.~522--526, 1991.

\bibitem{elgamal1985public}
T.~ElGamal,
``A Public Key Cryptosystem Based on Discrete Logarithms,''
\emph{IEEE Trans.\ Inf.\ Theory}, 31(4), 1985.

\bibitem{boneh2004short}
D.~Boneh, X.~Boyen, and H.~Shacham,
``Short Group Signatures,''
\emph{CRYPTO 2004}, pp.~41--55, 2004.

\bibitem{groth2016size}
J.~Groth,
``On the Size of Pairing-Based Non-interactive Arguments,''
\emph{EUROCRYPT 2016}, pp.~305--326, 2016.

\bibitem{diffie1976new}
W.~Diffie and M.~Hellman,
``New directions in cryptography,''
\emph{IEEE Trans.\ Inf.\ Theory}, 22(6), 1976.

\bibitem{marlinspike2016signal}
M.~Marlinspike and T.~Perrin,
``The Signal Protocol,''
2016.

\bibitem{danezis2009sphinx}
G.~Danezis and I.~Goldberg,
``Sphinx,''
\emph{IEEE S\&P}, 2009.

\bibitem{bernstein2006curve25519}
D.\,J.~Bernstein,
``Curve25519,''
\emph{PKC 2006}, 2006.

\bibitem{adida2008helios}
B.~Adida,
``Helios: Web-based Open-Audit Voting,''
\emph{USENIX Security}, pp.~335--348, 2008.

\bibitem{benaloh2006simple}
J.~Benaloh,
``Simple Verifiable Elections,''
\emph{EVT}, 2006.

\bibitem{amnesty2022iran}
Amnesty International,
``Iran: Digital Surveillance,''
2022.

\bibitem{reporters2023iran}
Reporters Without Borders,
``Iran: Digital Security for Journalists,''
2023.

\end{thebibliography}

\appendix

\section{Decoy Profiles}
\label{app:profiles}

\begin{table}[h]
\centering
\begin{tabular}{@{}lrrr@{}}
\toprule
\textbf{Profile} & \textbf{Balance} & \textbf{Tx/Wk} & \textbf{Tokens} \\
\midrule
Minimal  & 0.01--0.1 PARS & 1 & 0 \\
Standard & 0.1--2.0 PARS  & 3 & 2 \\
Active   & 1.0--10.0 PARS & 7 & 5 \\
\bottomrule
\end{tabular}
\end{table}

\section{Panic Sequence}
\label{app:sequence}

\begin{enumerate}
    \item Enter panic password
    \item Derive decoy key
    \item Background: rotate keys (100\,ms)
    \item Background: erase primary state (200\,ms)
    \item Background: mesh alerts (async)
    \item Foreground: present decoy ($<$50\,ms)
    \item Total visible delay: $<$300\,ms
\end{enumerate}

\end{document}
